\scanpage{42}

% B1 | text | opening paragraph and additions above lines
(Si on se donne la \deleted{\unsure{$\mathscr G$} \unsure{gerbe}} \inserted{catégorie de \unsure{Dessin}$_\otimes$ \unsure{polarisée}}
$(\unsure{C'|C'_f}(A'),T')$ sur $A'=\prod_{\ell\ne p,\infty}\mathbb Z_\ell$ comme
associée à la gerbe des racines carrées de $A'(1)$ (définie
par le choix d'une clôture algébrique $\bar{\mathbb F}_p$ de $\mathbb F_p$), \deleted{et} la
catégorie de Dessin$_\otimes$ $C(\mathbb Z_p)$ comme celle de $F$-cristaux sur $\bar{\mathbb F}_p$
(\unsure{par|pr} iso) qui sont de rang $2$, \deleted{\unsure{poids}} pente $\frac12$ \inserted{avec le $F$-cristal de Tate}, et la \deleted{\unsure{gerbe}}
catégorie $C(\mathbb Z_\infty)=C(\mathbb R)$ comme associée à la gerbe
type $\unsure{\mathscr G_\infty|\mathcal G_\infty}$ de groupe $\mu_2$ sur $\mathbb R$, on trouve en termes \deleted{\unsure{de}} $(C)$ un
certain torseur \inserted{$\mathcal T'(C)$} sous le groupe

% B2 | display | first product
\[
\deleted{\unsure{\mathbb Z_p}}\quad A'{}^*/A'{}^{*2}\times
(\mathbb Z_p^*/\mathbb Z_p^{*2}\times\mathbb Z/2\mathbb Z)
\times\mathbb R^*/\mathbb R^{*2}=\Gamma'
\]

% B3 | text | middle paragraph
Le groupe \inserted{$\mathbb Z/2\mathbb Z$} des classes d'isomorphie d'autoéquivalences de
\deleted{$C$} opère sur ce torseur, de façon compatible avec
son opération triviale sur lui-$=$, donc par translation,
et il opère donc via un hom \inserted{$f$ (\unsure{``Frobenius''})} de $\mathbb Z/2\mathbb Z$ dans ce groupe,
\deleted{\unsure{dont les}} l'image du \inserted{gén.} $\unsure{f}$ de $\mathbb Z/2\mathbb Z$ ds les 3 facteurs
étant respectivement

% B4 | display | middle lower, three components
\[
(p,\ \unsure{f|\beta}=(0,1),\ p)
\]

% B5 | text | small right-hand note linked to the last p
\margin{$\downarrow$ N.B. \deleted{\unsure{En fait là}, \unsure{pas} \unsure{légitime}},
\inserted{Comme $p\in\mathbb R^{*2}$} $\mathscr G_{\mathbb R}$ \unsure{orientable},
\unsure{entre} l'une des \unsure{ple}ssions,
\deleted{distinguer canoniqt \unsure{le}}
$\unsure{\mathrm{Pl}(C)}$ il y a un système transitif d'isomorphismes !}

% B6 | text | lower paragraph
Considérons le quotient

% B7 | display | lower centred quotient
\[
\Gamma=\Gamma'/f(\pm1)
\]

% B8 | text | last lines and addition above the last group
\deleted{On voit donc}
et soit $\mathcal T(C)$ le quotient de $\mathcal T'(C)$ par $\pm1$, qui
est donc un torseur sous $\Gamma$. Ce torseur, à isom unique
près, ne dépend pas de $C$, \inserted{on peut l'appeler $\mathcal T$ ou $\mathcal T(\bar{\mathbb F}_p)$.} D'ailleurs si $\Gamma'_0$ est le
ss-groupe de $\Gamma'$\quad $A'{}^*/A'{}^{*2}\times\mathbb Z_p^*/\mathbb Z_p^{*2}\times\mathbb R^*/\mathbb R^{*2}
=H^1_{\deleted{\unsure{\mathrm{ad}}}}(\mathbb A,\mu_2)$ \inserted{adèles}
on trouve que $\Gamma'\simeq\Gamma'_0\times f(\pm1)$, donc $\Gamma\simeq\Gamma'_0$ can, donc
% archive sheet number: 41
